\section{Introduction}
\label{sec:intro}

Self-supervised learning (SSL) methods have emerged as powerful techniques for learning transferable features across diverse tasks~\cite{chen2020simple, DINO, oquab2023dinov2}. However, their performance is significantly affected by dataset imbalance~\cite{assran2210hidden, tian2021divide, ts}. For instance, SimCLR underperforms on long-tailed datasets due to insufficient negative samples~\cite{tian2021divide}, while joint-embedding methods like VICReg assume uniform clustering, hampering performance on imbalanced data~\cite{assran2210hidden}.

Existing solutions typically adopt model-centric approaches, such as ensemble learning~\cite{tian2021divide}, incorporating arbitrary feature priors~\cite{assran2210hidden}, or modifying training dynamics~\cite{ts}. These approaches, however, often require prior knowledge of dataset distributions or extensive hyperparameter tuning, contradicting the unsupervised nature of SSL.

We propose BRADD (\textbf{B}alancing \textbf{R}epresentations with \textbf{A}nomaly \textbf{D}etection and \textbf{D}iffusion), a data-centric approach to address imbalance in SSL. BRADD divides training into cycles, identifies underrepresented samples via OOD detection after each cycle, and augments them using diffusion models for subsequent training. This approach (1) eliminates the need for prior dataset knowledge and (2) dynamically balances the latent space to avoid suboptimal local minima.

Experiments on ImageNet-100-LT across multiple architectures (ResNet-50, ViT-S, ViT-B) and SSL methods (SimCLR, DINO, MoCo) demonstrate that BRADD consistently outperforms both balanced and imbalanced baselines on diverse downstream tasks. BRADD achieves significant improvements on fine-grained tasks (up to 11.7\% on Oxford Flowers) and surpasses state-of-the-art methods on CIFAR-10 (73.3\%) and CIFAR-100 (45.4\%), showing that our data-centric approach offers a promising alternative to model-centric solutions.